\section{PANOPTES Positioning and Strategic Research Roadmap}
\label{sec:panoptes-positioning}

To bridge the architectural gaps synthesized in Table~\ref{tab:related_works} and operationalize the transition to \textbf{OMM Level~4 (Optimized)}, Figure~\ref{fig:research-roadmap} outlines a strategic research roadmap structured around three advancement pillars, explicitly positioning the contributions of the PANOPTES reference architecture within the broader 5-year cloud-native research horizon.

\begin{figure*}[!ht]
\centering
\resizebox{0.96\textwidth}{!}{%
\begin{tikzpicture}[
    node distance=0.4cm and 0.8cm,
    font=\sffamily\scriptsize,
    box_base/.style={rectangle, rounded corners=3pt, draw=black!60, thick, align=left, inner sep=5pt},
    gap_box/.style={box_base, fill=red!8, draw=red!60!black, text width=4.2cm, minimum height=1.45cm},
    pillar_box/.style={box_base, fill=cyan!8, draw=cyan!60!black, text width=5.2cm, minimum height=1.45cm},
    target_box/.style={box_base, fill=green!10, draw=green!50!black, text width=4.2cm, minimum height=1.45cm},
    header_box/.style={font=\sffamily\bfseries\footnotesize, align=center, fill=gray!20, draw=gray!60, rounded corners=2pt, inner sep=3.5pt},
    arrow/.style={-{Stealth[length=2.5mm, width=2mm]}, thick, draw=black!70}
]

    % Column Headers (Posicionamento relativo horizontal)
    \node[header_box, text width=4.4cm, minimum height=0.8cm] (h1) at (0, 0) {Empirical Gaps Identified (RQ04)};
    \node[header_box, text width=5.4cm, minimum height=0.8cm, right=0.8cm of h1] (h2) {Strategic Research Pillars (5-Year Horizon)};
    \node[header_box, text width=4.4cm, minimum height=0.8cm, right=0.8cm of h2] (h3) {Target Architectural Outcomes};

    % Row 1: Pillar 1
    \node[gap_box, below=0.3cm of h1] (g1) {
        \textbf{G1. Tooling Fragmentation \& Overhead}\\
        $\bullet$ Lossy OTLP--Prometheus translation \cite{hartikainen_2024}\\
        $\bullet$ High CPU/memory sidecar cost \cite{araujo_2024, zanatta_23}
    };
    \node[pillar_box, below=0.3cm of h2] (p1) {
        \textbf{Pillar 1: Semantic Interoperability \& Low-Overhead eBPF}\\
        $\bullet$ Harmonized OTLP/OpenMetrics schemas (histograms, labels)\\
        $\bullet$ Kernel-to-application context propagation via eBPF without code changes
    };
    \node[target_box, below=0.3cm of h3] (t1) {
        \textbf{Unified Telemetry Fabric}\\
        $\bullet$ $<2\%$ runtime CPU overhead\\
        $\bullet$ Vendor-neutral zero-code pipeline
    };

    % Row 2: Pillar 2
    \node[gap_box, below=0.3cm of g1] (g2) {
        \textbf{G2. Multimodal Correlation Silos}\\
        $\bullet$ High cardinality \& data overload \cite{deng_2024, han_2024}\\
        $\bullet$ Manual cross-pillar RCA queries \cite{gomes_2025, usman_2022}
    };
    \node[pillar_box, below=0.3cm of p1] (p2) {
        \textbf{Pillar 2: AIOps for Multimodal Telemetry Correlation}\\
        $\bullet$ Graph Attention Networks (GATs) over $O = \{M, L, T, E\}$ \cite{han_2024}\\
        $\bullet$ Topology-aware causal inference under pod ephemerality
    };
    \node[target_box, below=0.3cm of t1] (t2) {
        \textbf{Automated Root Cause Localization}\\
        $\bullet$ Drastic MTTR \& cognitive load reduction\\
        $\bullet$ Context-aware tail-based sampling
    };

    % Row 3: Pillar 3
    \node[gap_box, below=0.3cm of g2] (g3) {
        \textbf{G3. Reactive Control \& Blind Spots}\\
        $\bullet$ Decoupled security \& energy metrics \cite{deng_2024}\\
        $\bullet$ Static threshold autoscaling \cite{gomes_2025, kosinska_2023}
    };
    \node[pillar_box, below=0.3cm of p2] (p3) {
        \textbf{Pillar 3: Autonomic, Secure \& Sustainable Observability}\\
        $\bullet$ Closed-loop Kubernetes controllers driven by predictive SLOs\\
        $\bullet$ Convergence of runtime eBPF SecOps and carbon-aware GreenOps
    };
    \node[target_box, below=0.3cm of t2] (t3) {
        \textbf{OMM Level 5 (Autonomic Platform)}\\
        $\bullet$ Self-healing workload orchestration\\
        $\bullet$ Energy-proportional telemetry \& clusters
    };

    % Arrows connecting columns
    \draw[arrow] (g1.east) -- (p1.west);
    \draw[arrow] (p1.east) -- (t1.west);
    \draw[arrow] (g2.east) -- (p2.west);
    \draw[arrow] (p2.east) -- (t2.west);
    \draw[arrow] (g3.east) -- (p3.west);
    \draw[arrow] (p3.east) -- (t3.west);

\end{tikzpicture}%
}
\caption{Strategic Research Roadmap Mapping Empirical Gaps to Research Pillars and Target Architectural Outcomes.}
\label{fig:research-roadmap}
\end{figure*}

As illustrated in Figure~\ref{fig:research-roadmap}, the roadmap connects the empirical gaps identified in our Rapid Review (\textbf{G1}--\textbf{G3}) to three strategic pillars:

\begin{enumerate}
    \item \textbf{Pillar 1 -- Semantic Interoperability and Low-Overhead Instrumentation (Addressing G1):} Addresses tooling fragmentation and resource contention \cite{hartikainen_2024, araujo_2024, zanatta_23}. Within this thesis, PANOPTES directly instantiates the harmonized OTLP/Prometheus/Loki/Tempo pipeline using lightweight DaemonSet collection, label-only log indexing, and index-free object trace storage to keep dynamic CPU overhead within a strict operational envelope ($H_{1,3}$), while establishing the architectural interface for future zero-code eBPF kernel telemetry.
    
    \item \textbf{Pillar 2 -- Multimodal Telemetry Correlation and Automated RCA (Addressing G2):} Tackles multimodal correlation silos and cascading fault ambiguity \cite{deng_2024, han_2024, gomes_2025}. PANOPTES operationalizes this pillar at \textbf{OMM Level~4} through Control Principles \texttt{CP1} (\textit{Causal Linkage}) and \texttt{CP2} (\textit{Unified State Representation}), combining PromQL saturation/restart signals, LogQL exception streams, and Tempo exclusive span self-time decomposition ($T_{\text{self}} = T_{\text{span}} - \sum T_{\text{children}}$) with L4/L7 peer attribution to deterministically isolate root causes ($H_{1,1}$ and $H_{1,2}$).
    
    \item \textbf{Pillar 3 -- Autonomic, Secure, and Sustainable Observability (Addressing G3 -- Future Horizon toward OMM Level 5):} Targets the evolution from deterministic root-cause diagnosis (Level~4) to autonomous closed-loop actuation (Level~5), where the correlated telemetry graph produced by PANOPTES feeds predictive AIOps controllers, runtime eBPF SecOps enforcement, and carbon-aware GreenOps scheduling \cite{deng_2024, kosinska_2023}.
\end{enumerate}

By positioning PANOPTES as the concrete, reproducible open-source instantiation of \textbf{OMM Level~4}, this thesis resolves the foundational correlation and overhead bottlenecks (\textbf{G1} and \textbf{G2}) that must be solved before autonomous Level~5 self-healing systems can be reliably deployed in production.
